\documentclass[12pt,a4paper,twoside]{report}

%% Packages
\usepackage[french]{babel}
\usepackage[T1]{fontenc}
\usepackage{geometry}
\usepackage{graphicx}
\usepackage{listings}
\usepackage{xcolor}
\usepackage{fancyhdr}
\usepackage{setspace}
\usepackage{hyperref}
\usepackage{tocloft}
\usepackage{array}
\usepackage{booktabs}
\usepackage{float}
\usepackage{amsmath}
\usepackage{amssymb}
\usepackage{tikz}
\usepackage{caption}
\usepackage{subcaption}
\usepackage{pdfpages}

%% Geometry
\geometry{left=3cm, right=2cm, top=2.5cm, bottom=2.5cm}
\setstretch{1.5}

%% Colors
\definecolor{codeblue}{rgb}{0.13, 0.29, 0.53}
\definecolor{codegray}{rgb}{0.5, 0.5, 0.5}
\definecolor{codegreen}{rgb}{0.0, 0.5, 0.0}
\definecolor{codeviolet}{rgb}{0.73, 0.13, 0.93}
\definecolor{lightgray}{rgb}{0.95, 0.95, 0.95}
\definecolor{darkblue}{rgb}{0.0, 0.2, 0.5}

%% Listings Configuration
\lstset{
    basicstyle=\ttfamily\small,
    breaklines=true,
    columns=fullflexible,
    commentstyle=\color{codegreen},
    keywordstyle=\color{codeblue}\bfseries,
    numberstyle=\tiny\color{codegray},
    stringstyle=\color{codeviolet},
    backgroundcolor=\color{lightgray},
    tabsize=2,
    language=Python,
    numbers=left,
    numbersep=5pt,
    showstringspaces=false,
    breakatwhitespace=true,
    frame=single,
    rulecolor=\color{darkblue},
    framexleftmargin=15pt,
    xleftmargin=15pt,
}

%% Header and Footer
\pagestyle{fancy}
\fancyhf{}
\fancyhead[LE,RO]{\textit{Projet DevOps - Task Manager}}
\fancyhead[LO,RE]{MEIA Cloud DevOps Engineering}
\fancyfoot[C]{\thepage}
\renewcommand{\headrulewidth}{0.5pt}
\renewcommand{\footrulewidth}{0.5pt}

%% TOC Style
\renewcommand{\cftchapfont}{\bfseries\Large}
\renewcommand{\cftchappagefont}{\bfseries}
\renewcommand{\cftsecfont}{\large}
\renewcommand{\cftsubsecfont}{\normalsize}

%% Hyperlink Setup
\hypersetup{
    colorlinks=true,
    linkcolor=darkblue,
    filecolor=darkblue,
    urlcolor=darkblue,
    citecolor=darkblue,
}

%% Custom Title Page
\title{\vspace{2cm}\textbf{\Huge Rapport de Projet DevOps}\\\textbf{\Large Déploiement d'une API REST sur Azure avec CI/CD}\vspace{2cm}}
\author{\Large \textit{Auteurs :}\\[0.5cm]\Large LEGNIOUI SIFEDDINE\\LOUAT OUSSAMA\\AYOUB AFRAOUI\\KAWTAR ELBEJJAJI\\\vspace{2cm}\large Encadrant Académique\\Module: MEIA - Cloud DevOps Engineering}
\date{\large \today}

%% Begin Document
\begin{document}

%% Title Page
\maketitle
\thispagestyle{empty}
\newpage

%% Abstract
\begin{abstract}
\noindent
Ce rapport présente un projet complet d'ingénierie DevOps réalisé dans le cadre du module MEIA Cloud DevOps Engineering. Le projet "Task Manager API" est une démonstration pratique de l'ensemble de la chaîne DevOps moderne, incluant le développement d'une application, la conteneurisation, l'infrastructure as code, la configuration automatisée et l'orchestration de conteneurs.

L'objectif principal est de déployer une application Flask REST sur une infrastructure cloud (Azure) en utilisant les meilleures pratiques de l'industrie DevOps. Le projet intègre Docker pour la conteneurisation, Terraform pour la provision d'infrastructure, Ansible pour la configuration des serveurs, Kubernetes (K3s) pour l'orchestration des conteneurs, et GitHub Actions pour l'intégration et le déploiement continus.

Ce document fournit une documentation détaillée de tous les composants, des principaux livrables, de la stratégie de déploiement et des scripts utilisés. Il constitue un guide complet pour reproduire l'infrastructure et comprendre les choix architecturaux.

\vspace{1cm}
\textbf{Mots-clés :} DevOps, Cloud Computing, Azure, Docker, Kubernetes, Terraform, Ansible, CI/CD, Infrastructure as Code, Orchestration
\end{abstract}

\newpage

%% Table of Contents
\tableofcontents
\newpage

%% Chapter 1: Introduction
\chapter{Introduction}

\section{Contexte du projet}

Le projet "Task Manager API" est conçu comme une application d'apprentissage pratique pour les concepts avancés de DevOps et d'ingénierie cloud. Dans le contexte actuel de digitalisation et de transformations cloud, la maîtrise des technologies et des méthodologies DevOps est devenue essentielle pour tout ingénieur infrastructurePédagogique.

Ce projet couvre l'intégralité de la chaîne de valeur DevOps:
\begin{itemize}
    \item \textbf{Développement} : Une application REST API légère et modulaire
    \item \textbf{Containerisation} : Empaquetage de l'application dans des conteneurs Docker
    \item \textbf{Infrastructure} : Provisioning automatisé d'infrastructure sur Azure via Terraform
    \item \textbf{Configuration} : Automation de configuration serveur via Ansible
    \item \textbf{Orchestration} : Déploiement et gestion des conteneurs via Kubernetes (K3s)
    \item \textbf{CI/CD} : Pipeline automatisé via GitHub Actions
\end{itemize}

\section{Objectifs du projet}

Les objectifs principaux de ce projet sont:

\begin{enumerate}
    \item Développer une application REST complète avec Flask
    \item Containeriser l'application de manière optimale avec Docker
    \item Automatiser le provisioning d'infrastructure cloud (Azure) avec Terraform
    \item Configurer automatiquement les serveurs avec Ansible
    \item Orchestrer les conteneurs avec Kubernetes (K3s)
    \item Implémenter un pipeline CI/CD complet avec GitHub Actions
    \item Documenter l'ensemble de l'infrastructure et du processus de déploiement
    \item Démontrer les meilleures pratiques DevOps dans un contexte réel
\end{enumerate}

\section{Public cible et utilité du document}

Ce rapport s'adresse à:
\begin{itemize}
    \item Les étudiants en ingénierie cloud et DevOps
    \item Les administrateurs systèmes intéressés par l'automatisation
    \item Les ingénieurs DevOps souhaitant améliorer leurs compétences
    \item Toute personne intéressée par les architectures cloud modernes
\end{itemize}

Le document sert de guide complet pour comprendre, reproduire et améliorer la solution proposée.

\section{Structure du rapport}

\begin{description}
    \item[\textbf{Chapitre 2}] Fournit une vue globale de l'architecture et des composants
    \item[\textbf{Chapitre 3}] Détaille l'application Flask développée
    \item[\textbf{Chapitre 4}] Explique la containerisation Docker
    \item[\textbf{Chapitre 5}] Décrit l'infrastructure as code avec Terraform
    \item[\textbf{Chapitre 6}] Explique la configuration automatisée avec Ansible
    \item[\textbf{Chapitre 7}] Détaille l'orchestration Kubernetes
    \item[\textbf{Chapitre 8}] Présente le pipeline CI/CD
    \item[\textbf{Chapitre 9}] Fournit des instructions de déploiement
    \item[\textbf{Chapitre 10}] Couvre les aspects de maintenance et monitoring
    \item[\textbf{Annexes}] Code source complet et configurations
\end{description}

\newpage

%% Chapter 2: Architecture Générale
\chapter{Architecture Générale du Projet}

\section{Vue d'ensemble de l'architecture}

Le projet suit une architecture moderne et modulaire exploitant les principes de l'Infrastructure as Code, de la Conteneurisation et de l'Orchestration. La figure ci-dessous illustre les différentes couches de l'architecture.

\begin{figure}[H]
    \centering
    \caption{Architecture générale du projet DevOps}
    \begin{tikzpicture}[scale=0.9, every node/.style={draw, rectangle, text centered, minimum width=3cm}]
        \node[fill=lightgray] (github) at (0, 8) {GitHub Repository};
        \node[fill=lightgray] (github_actions) at (0, 6.5) {GitHub Actions\\CI/CD Pipeline};
        \node[fill=cyan!30] (docker_hub) at (0, 5) {Docker Hub\\Image Registry};
        
        \node[fill=yellow!30] (azure) at (3.5, 7) {Azure Cloud};
        \node[fill=orange!30] (terraform) at (3.5, 5.5) {Terraform\\Infrastructure};
        \node[fill=green!30] (vnet) at (3.5, 4) {Virtual Network\\+ NSG};
        \node[fill=red!30] (master) at (1.5, 2) {Master Node\\K3s + Docker};
        \node[fill=red!30] (worker) at (5.5, 2) {Worker Node\\K3s + Docker};
        \node[fill=purple!30] (ansible) at (3.5, 0.5) {Ansible\\Configuration};
        
        \node[fill=blue!20] (app) at (1.5, -1.5) {Task Manager\\Container};
        \node[fill=blue!20] (app2) at (5.5, -1.5) {Task Manager\\Container};
        
        \draw[->] (github) -- (github_actions);
        \draw[->] (github_actions) -- (docker_hub);
        \draw[->] (docker_hub) -- (azure);
        \draw[->] (terraform) -- (vnet);
        \draw[->] (vnet) -- (master);
        \draw[->] (vnet) -- (worker);
        \draw[->] (ansible) -- (master);
        \draw[->] (ansible) -- (worker);
        \draw[->] (docker_hub) -- (app);
        \draw[->] (docker_hub) -- (app2);
    \end{tikzpicture}
\end{figure}

\section{Composants principaux}

\subsection{Application Flask}
Une application REST légère écrite en Python utilisant le framework Flask. Elle expose des endpoints pour la gestion de tâches (CRUD operations).

\subsection{Conteneurisation Docker}
Docker est utilisé pour packager l'application avec toutes ses dépendances dans un conteneur isolé et reproductible. Cela assure la cohérence entre les environnements de développement, test et production.

\subsection{Infrastructure Azure}
L'infrastructure est hébergée sur Microsoft Azure, utilisant:
\begin{itemize}
    \item Resource Groups pour l'organisation logique
    \item Virtual Networks (VNet) pour la communication sécurisée
    \item Network Security Groups (NSG) pour le contrôle d'accès
    \item Virtual Machines (VMs) Linux Ubuntu pour les nœuds
\end{itemize}

\subsection{Terraform}
Infrastructure as Code (IaC) pour l'automatisation complète du provisioning Azure. Permet une réproductibilité et une gestion de version de l'infrastructure.

\subsection{Ansible}
Automation de configuration pour:
\begin{itemize}
    \item Installation de Docker sur chaque nœud
    \item Installation et configuration de K3s (Kubernetes léger)
    \item Jointure des workers au cluster master
\end{itemize}

\subsection{Kubernetes (K3s)}
Orchestration de conteneurs utilisant K3s, une distribution Kubernetes légère idéale pour les environnements edge et de test.

\subsection{Pipeline CI/CD}
GitHub Actions automatise:
\begin{itemize}
    \item Les tests et linting du code
    \item La construction de l'image Docker
    \item Le push vers Docker Hub
    \item Le déploiement sur le cluster Kubernetes
\end{itemize}

\section{Flux de déploiement}

\begin{enumerate}
    \item Les développeurs poussent le code vers GitHub
    \item GitHub Actions déclenche automatiquement le pipeline CI/CD
    \item L'application est testée et construite
    \item L'image Docker est créée et poussée vers Docker Hub
    \item Kubernetes récupère la nouvelle image
    \item Les pods sont mis à jour avec le rolling deployment
    \item L'application est accessible via le service NodePort
\end{enumerate}

\section{Avantages de cette architecture}

\begin{itemize}
    \item \textbf{Scalabilité} : Kubernetes permet de scaler l'application facilement
    \item \textbf{Automatisation} : Terraform et Ansible automatisent tout le provisioning
    \item \textbf{Reproductibilité} : Infrastructure as Code garantit la reproductibilité
    \item \textbf{Continuous Delivery} : Pipeline CI/CD pour des déploiements rapides
    \item \textbf{Résilience} : K3s assure la haute disponibilité
    \item \textbf{Isolation} : Docker isolate complètement l'application
    \item \textbf{Flexibilité} : Facile à adapter et à modifier
\end{itemize}

\newpage

%% Chapter 3: Application Flask
\chapter{Application Flask - Task Manager API}

\section{Présentation générale}

L'application Task Manager API est une API REST complète écrite en Python avec le framework Flask. Elle permet de gérer une liste de tâches avec les opérations CRUD (Create, Read, Update, Delete) standard.

\section{Architecture de l'application}

L'application suit une architecture simple et monolithique appropriée pour une démonstration DevOps. Elle utilise:
\begin{itemize}
    \item \textbf{Flask} : Framework web minimaliste et flexible
    \item \textbf{Python 3.9} : Langage de programmation moderne
    \item \textbf{Gunicorn} : Serveur WSGI pour la production
    \item \textbf{UUID} : Pour la génération d'identifiants uniques
\end{itemize}

\section{Endpoints API}

\begin{table}[H]
    \centering
    \caption{Endpoints de l'API Task Manager}
    \begin{tabular}{|l|l|l|p{4cm}|}
        \hline
        \textbf{Méthode} & \textbf{Endpoint} & \textbf{Description} & \textbf{Code} \\
        \hline
        GET & \texttt{/} & Retourne la page HTML & 200 \\
        GET & \texttt{/tasks} & Liste toutes les tâches & 200 \\
        POST & \texttt{/tasks} & Crée une nouvelle tâche & 201 \\
        PATCH/PUT & \texttt{/tasks/<id>} & Met à jour une tâche & 200/404 \\
        DELETE & \texttt{/tasks/<id>} & Supprime une tâche & 200/404 \\
        GET & \texttt{/health} & Health check & 200 \\
        \hline
    \end{tabular}
\end{table}

\section{Structure des données}

Une tâche est représentée par la structure JSON suivante:

\begin{lstlisting}[language=json]
{
    "id": "uuid-string",
    "title": "Titre de la tâche",
    "description": "Description optionnelle",
    "done": false
}
\end{lstlisting}

\section{Code source principal}

\subsection{Initialisation et routes principales}

\begin{lstlisting}[language=Python]
from flask import Flask, jsonify, request, render_template
import uuid

app = Flask(__name__)
tasks = []

@app.route('/')
def index():
    return render_template('index.html')

@app.route('/health', methods=['GET'])
def health_check():
    return jsonify({'status': 'healthy'}), 200
\end{lstlisting}

\subsection{Opération GET - Récupérer toutes les tâches}

\begin{lstlisting}[language=Python]
@app.route('/tasks', methods=['GET'])
def get_tasks():
    return jsonify(tasks), 200
\end{lstlisting}

Cette route retourne toutes les tâches actuellement stockées dans la liste en mémoire.

\subsection{Opération POST - Créer une tâche}

\begin{lstlisting}[language=Python]
@app.route('/tasks', methods=['POST'])
def create_task():
    data = request.get_json()
    if not data or not 'title' in data:
        return jsonify({'error': 'Title is required'}), 400
    
    task = {
        'id': str(uuid.uuid4()),
        'title': data['title'],
        'description': data.get('description', ''),
        'done': False
    }
    tasks.append(task)
    return jsonify(task), 201
\end{lstlisting}

\subsection{Opération PATCH/PUT - Mettre à jour une tâche}

\begin{lstlisting}[language=Python]
@app.route('/tasks/<task_id>', methods=['PATCH', 'PUT'])
def update_task(task_id):
    data = request.get_json()
    for task in tasks:
        if task['id'] == task_id:
            if 'done' in data:
                task['done'] = data['done']
            return jsonify(task), 200
    return jsonify({'error': 'Task not found'}), 404
\end{lstlisting}

\subsection{Opération DELETE - Supprimer une tâche}

\begin{lstlisting}[language=Python]
@app.route('/tasks/<task_id>', methods=['DELETE'])
def delete_task(task_id):
    global tasks
    tasks = [task for task in tasks if task['id'] != task_id]
    return jsonify({'result': True}), 200
\end{lstlisting}

\section{Interface utilisateur}

L'application inclut une interface web moderne avec:
\begin{itemize}
    \item Design responsive utilisant CSS moderne
    \item Upload de tâches via formulaire
    \item Affichage dynamique via JavaScript
    \item Icones Font Awesome pour une meilleure UX
\end{itemize}

\section{Dépendances Python}

\begin{lstlisting}
Flask==3.0.0
Werkzeug==3.0.1
gunicorn==21.2.0
\end{lstlisting}

\section{Démarrage de l'application}

En développement (avec Flask):
\begin{lstlisting}[language=bash]
python app.py
\end{lstlisting}

En production (avec Gunicorn):
\begin{lstlisting}[language=bash]
gunicorn --bind 0.0.0.0:5000 app:app
\end{lstlisting}

\newpage

%% Chapter 4: Containerisation Docker
\chapter{Containerisation Docker}

\section{Vue d'ensemble de Docker}

Docker est une plateforme de containerisation qui permet d'empaqueter une application avec toutes ses dépendances dans un conteneur isolé. Les avantages principaux sont:

\begin{itemize}
    \item \textbf{Reproducibilité} : Même comportement sur tous les environnements
    \item \textbf{Isolation} : Application isolée du système hôte
    \item \textbf{Légèreté} : Plus léger qu'une machine virtuelle
    \item \textbf{Portabilité} : Fonctionne sur tout système avec Docker
    \item \textbf{Scalabilité} : Facile de scaler avec l'orchestration
\end{itemize}

\section{Dockerfile}

Le Dockerfile définit comment construire l'image Docker de notre application.

\subsection{Analyse du Dockerfile}

\begin{lstlisting}[language=dockerfile]
# Base image
FROM python:3.9-slim

# Répertoire de travail
WORKDIR /app

# Copie des requirements
COPY app/requirements.txt .

# Installation des dépendances
RUN pip install --no-cache-dir -r requirements.txt

# Copie du code source
COPY app/ .

# Port exposé
EXPOSE 5000

# Commande de démarrage
CMD ["gunicorn", "--bind", "0.0.0.0:5000", "app:app"]
\end{lstlisting}

\subsection{Explications détaillées}

\begin{description}
    \item[\textbf{FROM}] Utilise l'image officielle Python 3.9-slim comme base (image minimale pour réduire la taille)
    \item[\textbf{WORKDIR}] Définit le répertoire de travail à \texttt{/app} dans le conteneur
    \item[\textbf{COPY requirements.txt}] Copie le fichier des dépendances
    \item[\textbf{RUN pip install}] Installe les dépendances Python (--no-cache-dir réduit la taille)
    \item[\textbf{COPY app/}] Copie le code source complet
    \item[\textbf{EXPOSE}] Déclare le port 5000 (informatif, ne fait pas d'exposition réelle)
    \item[\textbf{CMD}] Commande à exécuter au démarrage du conteneur
\end{description}

\section{Meilleure pratiques appliquées}

\begin{enumerate}
    \item \textbf{Image de base minimaliste} : Utilisation de python:3.9-slim au lieu de python:3.9
    \item \textbf{Layer caching} : Dépendances copiées avant le code pour meilleur caching
    \item \textbf{Réduction de la taille} : --no-cache-dir lors de pip install
    \item \textbf{Utilisateur non-root} : (Peut être amélioré)
    \item \textbf{Serveur de production} : Gunicorn plutôt que Flask en développement
    \item \textbf{Single responsibility} : Un conteneur pour une seule application
\end{enumerate}

\section{Construction de l'image}

\subsection{Commande de base}

\begin{lstlisting}[language=bash]
docker build -t task-manager:latest .
\end{lstlisting}

\subsection{Construction avec tag complet}

\begin{lstlisting}[language=bash]
docker build -t seiflegnioui/task-manager:latest .
\end{lstlisting}

\subsection{Avec tag de version}

\begin{lstlisting}[language=bash]
docker build -t seiflegnioui/task-manager:v1.0.0 .
\end{lstlisting}

\section{Exécution du conteneur}

\subsection{Environnement de développement}

\begin{lstlisting}[language=bash]
docker run -p 5000:5000 task-manager:latest
\end{lstlisting}

\subsection{Environnement de production}

\begin{lstlisting}[language=bash]
docker run -d \
  --name task-manager \
  -p 80:5000 \
  --restart unless-stopped \
  task-manager:latest
\end{lstlisting}

\section{Optimisation des images Docker}

\subsection{Considérations de taille}

\begin{table}[H]
    \centering
    \caption{Taille estimée de l'image Docker}
    \begin{tabular}{|l|r|}
        \hline
        \textbf{Composant} & \textbf{Taille} \\
        \hline
        Base Python 3.9-slim & $\sim$ 150 MB \\
        Dépendances Python & $\sim$ 50 MB \\
        Code application & $\sim$ 1 MB \\
        \hline
        \textbf{Total} & $\sim$ 200 MB \\
        \hline
    \end{tabular}
\end{table}

\subsection{Multi-stage builds (recommandé)}

Une approche avancée utiliserait un multi-stage build:

\begin{lstlisting}[language=dockerfile]
# Stage 1: Builder
FROM python:3.9-slim as builder
WORKDIR /app
COPY app/requirements.txt .
RUN pip install --user --no-cache-dir -r requirements.txt

# Stage 2: Runtime
FROM python:3.9-slim
WORKDIR /app
COPY --from=builder /root/.local /root/.local
COPY app/ .
ENV PATH=/root/.local/bin:$PATH
CMD ["gunicorn", "--bind", "0.0.0.0:5000", "app:app"]
\end{lstlisting}

\section{Registre d'images}

Les images sont stockées dans Docker Hub pour faciliter le déploiement:

\begin{lstlisting}[language=bash]
# Login
docker login -u username

# Push
docker push seiflegnioui/task-manager:latest

# Pull
docker pull seiflegnioui/task-manager:latest
\end{lstlisting}

\newpage

%% Chapter 5: Infrastructure as Code avec Terraform
\chapter{Infrastructure as Code avec Terraform}

\section{Introduction à Terraform}

Terraform est un outil de Infrastructure as Code (IaC) qui permet de décrire et de gérer l'infrastructure cloud en tant que code. Les avantages principaux sont:

\begin{itemize}
    \item \textbf{Déclaratif} : Description de l'état désiré
    \item \textbf{Versionnable} : Infrastructure versionnée comme du code
    \item \textbf{Réutilisable} : Modules réutilisables
    \item \textbf{Prédictif} : Aperçu des changements avant application
    \item \textbf{Multi-cloud} : Support de plusieurs fournisseurs
\end{itemize}

\section{Configuration des providers}

\subsection{providers.tf}

\begin{lstlisting}[language=terraform]
terraform {
  required_version = ">= 1.0.0"

  required_providers {
    azurerm = {
      source  = "hashicorp/azurerm"
      version = "~> 3.0"
    }
    tls = {
      source  = "hashicorp/tls"
      version = "~> 4.0"
    }
    local = {
      source  = "hashicorp/local"
      version = "~> 2.1"
    }
  }
}

provider "azurerm" {
  features {}
}
\end{lstlisting}

\subsection{Explication des providers}

\begin{description}
    \item[\textbf{azurerm}] Provider Microsoft Azure pour la gestion des ressources Azure
    \item[\textbf{tls}] Provider pour la génération de clés et certificats TLS/SSL
    \item[\textbf{local}] Provider pour la gestion des fichiers locaux (inventory Ansible)
\end{description}

\section{Variables Terraform}

\subsection{variables.tf}

\begin{lstlisting}[language=terraform]
variable "resource_group_name" {
  type        = string
  description = "Name of the resource group"
  default     = "rg-devops-project"
}

variable "location" {
  type        = string
  description = "Azure region for the resources"
  default     = "swedencentral"
}

variable "admin_username" {
  type        = string
  description = "Admin username for the VMs"
  default     = "azureuser"
}

variable "vm_size_master" {
  type        = string
  description = "Size of the Master VM"
  default     = "Standard_D2s_v3"
}

variable "vm_size_worker" {
  type        = string
  description = "Size of the Worker VM"
  default     = "Standard_D2s_v3"
}
\end{lstlisting}

\section{Ressources principales}

\subsection{Resource Group}

\begin{lstlisting}[language=terraform]
resource "azurerm_resource_group" "rg" {
  name     = var.resource_group_name
  location = var.location
}
\end{lstlisting}

\subsection{Virtual Network}

\begin{lstlisting}[language=terraform]
resource "azurerm_virtual_network" "vnet" {
  name                = "vnet-devops"
  address_space       = ["10.0.0.0/16"]
  location            = azurerm_resource_group.rg.location
  resource_group_name = azurerm_resource_group.rg.name
}

resource "azurerm_subnet" "subnet" {
  name                 = "subnet-devops"
  resource_group_name  = azurerm_resource_group.rg.name
  virtual_network_name = azurerm_virtual_network.vnet.name
  address_prefixes     = ["10.0.1.0/24"]

  depends_on = [
    azurerm_virtual_network.vnet
  ]
}
\end{lstlisting}

\subsection{Network Security Group}

\begin{lstlisting}[language=terraform]
resource "azurerm_network_security_group" "nsg" {
  name                = "nsg-devops"
  location            = azurerm_resource_group.rg.location
  resource_group_name = azurerm_resource_group.rg.name

  security_rule {
    name                       = "SSH"
    priority                   = 1001
    direction                  = "Inbound"
    access                     = "Allow"
    protocol                   = "Tcp"
    source_port_range          = "*"
    destination_port_range     = "22"
    source_address_prefix      = "*"
    destination_address_prefix = "*"
  }

  security_rule {
    name                       = "KubeAPI"
    priority                   = 1002
    direction                  = "Inbound"
    access                     = "Allow"
    protocol                   = "Tcp"
    source_port_range          = "*"
    destination_port_range     = "6443"
    source_address_prefix      = "*"
    destination_address_prefix = "*"
  }

  security_rule {
    name                       = "HTTP"
    priority                   = 1003
    direction                  = "Inbound"
    access                     = "Allow"
    protocol                   = "Tcp"
    source_port_range          = "*"
    destination_port_range     = "30080"
    source_address_prefix      = "*"
    destination_address_prefix = "*"
  }
}
\end{lstlisting}

\section{Outputs}

\subsection{outputs.tf}

\begin{lstlisting}[language=terraform]
data "azurerm_public_ip" "master" {
  name                = azurerm_public_ip.master_ip.name
  resource_group_name = azurerm_linux_virtual_machine.master.resource_group_name
}

data "azurerm_public_ip" "worker" {
  name                = azurerm_public_ip.worker_ip.name
  resource_group_name = azurerm_linux_virtual_machine.worker.resource_group_name
}

output "master_public_ip" {
  value = data.azurerm_public_ip.master.ip_address
}

output "worker_public_ip" {
  value = data.azurerm_public_ip.worker.ip_address
}

# Generate Ansible Inventory dynamically
resource "local_file" "ansible_inventory" {
  content = templatefile("${path.module}/inventory.ini.tpl", {
    master_ip = data.azurerm_public_ip.master.ip_address
    worker_ip = data.azurerm_public_ip.worker.ip_address
    user      = var.admin_username
  })
  filename = "${path.module}/../ansible/inventory.ini"
}
\end{lstlisting}

\section{Workflow de déploiement}

\subsection{Initialisation}

\begin{lstlisting}[language=bash]
cd terraform
terraform init
\end{lstlisting}

Cette commande:
\begin{itemize}
    \item Télécharge les providers requis
    \item Initialise l'état Terraform
    \item Prépare le répertoire pour les opérations
\end{itemize}

\subsection{Plan}

\begin{lstlisting}[language=bash]
terraform plan
\end{lstlisting}

Affiche les changements qui seront appliqués.

\subsection{Apply}

\begin{lstlisting}[language=bash]
terraform apply --auto-approve
\end{lstlisting}

Crée les ressources Azure.

\subsection{Destruction}

\begin{lstlisting}[language=bash]
terraform destroy --auto-approve
\end{lstlisting}

Supprime toutes les ressources créées.

\section{État Terraform}

\subsection{Fichiers d'état}

\begin{itemize}
    \item \textbf{terraform.tfstate} : État courant
    \item \textbf{terraform.tfstate.backup} : Sauvegarde de l'état précédent
\end{itemize}

\subsection{Bonnes pratiques}

\begin{itemize}
    \item Toujours faire des backups de l'état
    \item Stocker l'état dans un backend sécurisé (Azure Storage)
    \item Ne jamais committer l'état dans Git
    \item Utiliser des locks pour les opérations concurrentes
\end{itemize}

\newpage

%% Chapter 6: Configuration avec Ansible
\chapter{Configuration Automatisée avec Ansible}

\section{Introduction à Ansible}

Ansible est un outil de configuration management et automatisation de l'IT. Pour ce projet:
\begin{itemize}
    \item Installation de Docker sur tous les nœuds
    \item Installation de K3s sur le nœud master
    \item Jointure des workers au cluster master
    \item Récupération du kubeconfig
\end{itemize}

\section{Playbook principal}

\subsection{playbook.yml}

\begin{lstlisting}[language=yaml]
---
- name: Setup K3s Master
  hosts: master
  become: yes
  tasks:
    - name: Install Docker
      apt:
        name: docker.io
        state: present
        update_cache: yes

    - name: Download and install K3s on Master
      shell: curl -sfL https://get.k3s.io | sh -s - server --docker

    - name: Get K3s token
      slurp:
        src: /var/lib/rancher/k3s/server/node-token
      register: k3s_token

    - name: Fetch Kubeconfig to local machine
      fetch:
        src: /etc/rancher/k3s/k3s.yaml
        dest: ../kubeconfig.yaml
        flat: yes

- name: Setup K3s Worker
  hosts: worker
  become: yes
  tasks:
    - name: Install Docker
      apt:
        name: docker.io
        state: present
        update_cache: yes

    - name: Join K3s Node to Cluster
      shell: "curl -sfL https://get.k3s.io | K3S_URL=https://{{ hostvars[groups['master'][0]]['inventory_hostname'] }}:6443 K3S_TOKEN={{ hostvars[groups['master'][0]]['k3s_token']['content'] | b64decode | trim }} sh -s - --docker"
\end{lstlisting}

\subsection{Structure et explication}

Le playbook est organisé en deux play:

\subsubsection{Play 1: Setup K3s Master}

\begin{enumerate}
    \item \textbf{Install Docker} : Installe Docker via apt
    \item \textbf{Install K3s} : Télécharge et installe K3s en mode serveur
    \item \textbf{Get Token} : Récupère le token d'authentification du cluster
    \item \textbf{Fetch Kubeconfig} : Récupère le fichier kubeconfig pour la machine locale
\end{enumerate}

\subsubsection{Play 2: Setup K3s Worker}

\begin{enumerate}
    \item \textbf{Install Docker} : Installe Docker via apt
    \item \textbf{Join Cluster} : Rejoint le cluster master en utilisant l'URL et le token
\end{enumerate}

\section{Fichier d'inventaire}

\subsection{inventory.ini.tpl}

\begin{lstlisting}
[master]
${master_ip} ansible_user=${user} ansible_ssh_private_key_file=../terraform/id_rsa_devops ansible_ssh_common_args='-o StrictHostKeyChecking=no'

[worker]
${worker_ip} ansible_user=${user} ansible_ssh_private_key_file=../terraform/id_rsa_devops ansible_ssh_common_args='-o StrictHostKeyChecking=no'

[k3s_cluster:children]
master
worker
\end{lstlisting}

\section{Exécution du playbook}

\subsection{Prérequis}

\begin{itemize}
    \item Ansible installé localement
    \item SSH configuré avec la clé privée
    \item Inventory.ini généré par Terraform
\end{itemize}

\subsection{Commandes}

\subsubsection{Première exécution avec vérification SSH}

\begin{lstlisting}[language=bash]
cd ansible
chmod 600 ../terraform/id_rsa_devops
export ANSIBLE_HOST_KEY_CHECKING=False
ansible-playbook -i inventory.ini playbook.yml
\end{lstlisting}

\subsubsection{Exécutions ultérieures}

\begin{lstlisting}[language=bash]
ansible-playbook -i inventory.ini playbook.yml
\end{lstlisting}

\subsubsection{Mode verbose pour debug}

\begin{lstlisting}[language=bash]
ansible-playbook -i inventory.ini playbook.yml -vvv
\end{lstlisting}

\section{Variables et fact gathering}

\subsection{Utilisation de variables}

\begin{itemize}
    \item \textbf{hostvars} : Accès aux variables de hosts
    \item \textbf{groups} : Accès aux groupes d'inventory
    \item \textbf{b64decode} : Décodage de base64 pour le token
\end{itemize}

\subsection{Enregistrement de variables}

\begin{lstlisting}[language=yaml]
- name: Get K3s token
  slurp:
    src: /var/lib/rancher/k3s/server/node-token
  register: k3s_token
\end{lstlisting}

La variable \texttt{k3s\_token} peut être utilisée par les plays suivants.

\section{Bonnes pratiques}

\begin{itemize}
    \item Utilise \texttt{become: yes} pour l'escalade de privilèges
    \item Diable la vérification SSH pour la première connexion
    \item Utilise \texttt{slurp} pour lire les fichiers distants
    \item Utilise \texttt{fetch} pour télécharger des fichiers
    \item Organise les tâches de manière logique et lisible
\end{itemize}

\newpage

%% Chapter 7: Orchestration Kubernetes
\chapter{Orchestration Kubernetes (K3s)}

\section{Vue d'ensemble de Kubernetes}

Kubernetes est une plateforme d'orchestration de conteneurs qui automatise le déploiement, la mise à l'échelle et la gestion des applications conteneurisées. Pour ce projet, nous utilisons K3s (Kubernetes léger).

\subsection{Concepts clés}

\begin{description}
    \item[\textbf{Pod}] Unité de déploiement la plus petite, contenant un ou plusieurs conteneurs
    \item[\textbf{Deployment}] Contrôleur Kubernetes gérant les réplicas de pods
    \item[\textbf{Service}] Exposant les pods en tant que service réseau
    \item[\textbf{Node}] Machine (VM) exécutant les pods
    \item[\textbf{Cluster}] Ensemble de nodes Kubernetes
    \item[\textbf{kubelet}] Agent exécutant les pods sur chaque node
    \item[\textbf{API Server}] Interface de contrôle du cluster
\end{description}

\section{Manifests Kubernetes}

\subsection{Deployment}

\subsubsection{deployment.yaml}

\begin{lstlisting}[language=yaml]
apiVersion: apps/v1
kind: Deployment
metadata:
  name: task-manager-api
  labels:
    app: task-manager
spec:
  replicas: 2
  selector:
    matchLabels:
      app: task-manager
  template:
    metadata:
      labels:
        app: task-manager
    spec:
      containers:
      - name: task-manager-api
        image: seiflegnioui/task-manager:latest
        imagePullPolicy: Always
        ports:
        - containerPort: 5000
        resources:
          requests:
            memory: "128Mi"
            cpu: "100m"
          limits:
            memory: "256Mi"
            cpu: "250m"
        readinessProbe:
          httpGet:
            path: /health
            port: 5000
          initialDelaySeconds: 10
          periodSeconds: 5
\end{lstlisting}

\subsubsection{Explications détaillées}

\begin{description}
    \item[\textbf{apiVersion}] Version de l'API Kubernetes utilisée (apps/v1 pour Deployments)
    \item[\textbf{kind}] Type de ressource Kubernetes (Deployment)
    \item[\textbf{metadata}] Informations sur le déploiement (name, labels)
    \item[\textbf{spec.replicas}] Nombre d'instances de l'application (2)
    \item[\textbf{selector}] Sélecteur pour les pods gérés par ce déploiement
    \item[\textbf{template}] Template de pod définissant la configuration
    \item[\textbf{containers[].image}] Image Docker à déployer
    \item[\textbf{imagePullPolicy}] Always = toujours tirer la dernière image
    \item[\textbf{resources.requests}] Ressources minimales garanties
    \item[\textbf{resources.limits}] Ressources maximales autorisées
    \item[\textbf{readinessProbe}] Sonde de santé pour vérifier que l'app est prête
\end{description}

\subsection{Service}

\subsubsection{service.yaml}

\begin{lstlisting}[language=yaml]
apiVersion: v1
kind: Service
metadata:
  name: task-manager-service
spec:
  type: NodePort
  selector:
    app: task-manager
  ports:
    - protocol: TCP
      port: 80
      targetPort: 5000
      nodePort: 30080
\end{lstlisting}

\subsubsection{Types de service}

\begin{table}[H]
    \centering
    \caption{Types de service Kubernetes}
    \begin{tabular}{|l|p{5cm}|l|}
        \hline
        \textbf{Type} & \textbf{Description} & \textbf{Usage} \\
        \hline
        ClusterIP & Service interne au cluster & Communication interne \\
        NodePort & Expose sur tous les nodes & Accès temps de test \\
        LoadBalancer & Load balancer cloud & Production \\
        ExternalName & Alias DNS externe & Services externes \\
        \hline
    \end{tabular}
\end{table}

Pour ce projet, nous utilisons \textbf{NodePort} qui expose le service sur le port 30080 de chaque node.

\section{Déploiement sur Kubernetes}

\subsection{Prérequis}

\begin{itemize}
    \item Cluster K3s fonctionnel (installé par Ansible)
    \item \texttt{kubectl} configuré avec kubeconfig.yaml
    \item Image Docker disponible dans Docker Hub
\end{itemize}

\subsection{Commandes de déploiement}

\subsubsection{Appliquer les manifests}

\begin{lstlisting}[language=bash]
# Depuis le répertoire racine
kubectl apply -f kubernetes/

# Ou individuellement
kubectl apply -f kubernetes/deployment.yaml
kubectl apply -f kubernetes/service.yaml
\end{lstlisting}

\subsubsection{Vérifier le déploiement}

\begin{lstlisting}[language=bash]
# Lister les deployments
kubectl get deployments

# Lister les pods
kubectl get pods

# Lister les services
kubectl get services

# Détails d'un deployment
kubectl describe deployment task-manager-api

# Logs d'un pod
kubectl logs <pod-name>

# Logs en temps réel
kubectl logs -f <pod-name>
\end{lstlisting}

\subsection{Accès à l'application}

L'application est accessible via:
\begin{itemize}
    \item \textbf{URL} : \texttt{http://<master-ip>:30080}
    \item \textbf{API} : \texttt{http://<master-ip>:30080/tasks}
\end{itemize}

\section{Gestion du cycle de vie des pods}

\subsection{Mise à jour d'image}

\begin{lstlisting}[language=bash]
# Mise à jour du déploiement
kubectl set image deployment/task-manager-api \
  task-manager-api=seiflegnioui/task-manager:v2.0.0

# Ou avec kubectl apply après modification du YAML
kubectl apply -f kubernetes/deployment.yaml
\end{lstlisting}

\subsection{Rolling deployment}

Kubernetes effectue un rolling deployment par défaut:
\begin{itemize}
    \item Les anciens pods sont arrêtés progressivement
    \item Les nouveaux pods sont créés progressivement
    \item Zéro downtime garanti
\end{itemize}

\subsection{Scaling}

\begin{lstlisting}[language=bash]
# Scaler à 5 replicas
kubectl scale deployment task-manager-api --replicas=5

# Ou modifier le YAML
kubectl set env deployment/task-manager-api REPLICAS=5
\end{lstlisting}

\section{Health checks et probes}

\subsection{Readiness Probe}

Utilisé pour déterminer si le pod est prêt à recevoir du trafic:

\begin{lstlisting}[language=yaml]
readinessProbe:
  httpGet:
    path: /health
    port: 5000
  initialDelaySeconds: 10
  periodSeconds: 5
\end{lstlisting}

\subsection{Liveness Probe (recommandé)}

\begin{lstlisting}[language=yaml]
livenessProbe:
  httpGet:
    path: /health
    port: 5000
  initialDelaySeconds: 30
  periodSeconds: 10
  failureThreshold: 3
\end{lstlisting}

La liveness probe redémarre le pod si l'application ne répond pas.

\newpage

%% Chapter 8: Pipeline CI/CD
\chapter{Pipeline CI/CD avec GitHub Actions}

\section{Vue d'ensemble}

Le pipeline CI/CD automatise les étapes suivantes:
\begin{enumerate}
    \item Linting et tests du code
    \item Construction de l'image Docker
    \item Push vers Docker Hub
    \item Déploiement sur le cluster Kubernetes
\end{enumerate}

\section{Workflows GitHub Actions}

\subsection{Structure générale}

Une pipeline GitHub Actions est définie dans \texttt{.github/workflows/} avec les composants:

\begin{description}
    \item[\textbf{name}] Nom du workflow
    \item[\textbf{on}] Événements déclencheurs
    \item[\textbf{jobs}] Travaux à exécuter
    \item[\textbf{steps}] Étapes individuelles
\end{description}

\subsection{Exemple de workflow CI/CD}

\begin{lstlisting}[language=yaml]
name: Build and Deploy

on:
  push:
    branches: [ main ]
  pull_request:
    branches: [ main ]

jobs:
  build:
    runs-on: ubuntu-latest

    steps:
    - uses: actions/checkout@v3

    - name: Set up Python
      uses: actions/setup-python@v4
      with:
        python-version: 3.9

    - name: Install dependencies
      run: |
        python -m pip install --upgrade pip
        pip install -r app/requirements.txt
        pip install pytest pylint

    - name: Lint with pylint
      run: |
        pylint app/app.py --disable=C0111,W0612 || true

    - name: Test with pytest
      run: |
        pytest app/ -v || true

  deploy:
    needs: build
    runs-on: ubuntu-latest
    if: github.ref == 'refs/heads/main'

    steps:
    - uses: actions/checkout@v3

    - name: Set up Docker
      uses: docker/setup-buildx-action@v2

    - name: Login to Docker Hub
      uses: docker/login-action@v2
      with:
        username: \${{ secrets.DOCKER_USERNAME }}
        password: \${{ secrets.DOCKER_PASSWORD }}

    - name: Build and push
      uses: docker/build-push-action@v4
      with:
        context: .
        push: true
        tags: |
          \${{ secrets.DOCKER_USERNAME }}/task-manager:latest
          \${{ secrets.DOCKER_USERNAME }}/task-manager:v\${{ github.run_number }}

    - name: Deploy to Kubernetes
      run: |
        mkdir -p \$HOME/.kube
        echo "\${{ secrets.KUBECONFIG }}" | base64 -d > \$HOME/.kube/config
        kubectl apply -f kubernetes/
        kubectl rollout restart deployment/task-manager-api
\end{lstlisting}

\section{Déclencheurs (triggers)}

\subsection{Push sur la branche main}

\begin{lstlisting}[language=yaml]
on:
  push:
    branches: [ main ]
\end{lstlisting}

Le workflow s'exécute automatiquement à chaque push.

\subsection{Pull requests}

\begin{lstlisting}[language=yaml]
on:
  pull_request:
    branches: [ main ]
\end{lstlisting}

Le workflow s'exécute sur chaque PR.

\subsection{Schedules}

\begin{lstlisting}[language=yaml]
on:
  schedule:
    - cron: '0 2 * * *'  # Quotidien à 2h du matin
\end{lstlisting}

\section{Jobs et dépendances}

\subsection{Définition de jobs}

\begin{lstlisting}[language=yaml]
jobs:
  build:
    runs-on: ubuntu-latest
    # ... steps ...

  deploy:
    needs: build  # Dépens de build
    runs-on: ubuntu-latest
    # ... steps ...
\end{lstlisting}

\subsection{Conditions}

\begin{lstlisting}[language=yaml]
deploy:
  needs: build
  runs-on: ubuntu-latest
  if: github.ref == 'refs/heads/main'  # Seulement sur main
\end{lstlisting}

\section{Secrets et variables}

\subsection{Déclaration des secrets}

Dans GitHub: Settings $\rightarrow$ Secrets and variables $\rightarrow$ Actions

Secrets nécessaires:
\begin{itemize}
    \item \texttt{DOCKER\_USERNAME} : Username Docker Hub
    \item \texttt{DOCKER\_PASSWORD} : Token Docker Hub
    \item \texttt{KUBECONFIG} : Fichier kubeconfig en base64
\end{itemize}

\subsection{Utilisation dans le workflow}

\begin{lstlisting}[language=yaml]
- name: Login to Docker Hub
  uses: docker/login-action@v2
  with:
    username: \${{ secrets.DOCKER_USERNAME }}
    password: \${{ secrets.DOCKER_PASSWORD }}
\end{lstlisting}

\section{Actions réutilisables}

\subsection{Actions officielles}

\begin{itemize}
    \item \texttt{actions/checkout} : Cloner le repo
    \item \texttt{actions/setup-python} : Installer Python
    \item \texttt{docker/login-action} : Login Docker
    \item \texttt{docker/build-push-action} : Build et push
\end{itemize}

\subsection{Exemple: actions/checkout}

\begin{lstlisting}[language=yaml]
- uses: actions/checkout@v3
  with:
    fetch-depth: 0  # Fetch tous les commits
\end{lstlisting}

\section{Failure handling}

\subsection{Continue on error}

\begin{lstlisting}[language=yaml]
- name: Lint
  run: pylint app/
  continue-on-error: true  # Ne pas arrêter si erreur
\end{lstlisting}

\subsection{Conditions avancées}

\begin{lstlisting}[language=yaml]
- name: Step conditionnelle
  if: success()  # S'il n'y a pas eu d'erreur
  run: echo "Success!"

- name: Autre step conditionnelle
  if: failure()  # S'il y a eu une erreur
  run: echo "Failed!"

- name: Always run
  if: always()  # Toujours exécuter
  run: echo "Done"
\end{lstlisting}

\section{Monitoring et logging}

\subsection{Annotations}

\begin{lstlisting}[language=bash]
echo "::notice file=app.py,line=10::Found issue"
echo "::warning file=app.py,line=20::Be careful"
echo "::error file=app.py,line=30::Critical error"
\end{lstlisting}

\subsection{Groupes de logs}

\begin{lstlisting}[language=bash]
echo "::group::Titre du groupe"
# ... commandes ...
echo "::endgroup::"
\end{lstlisting}

\newpage

%% Chapter 9: Guide de Déploiement
\chapter{Guide Complet de Déploiement}

\section{Prérequis préalables}

\subsection{Logiciels à installer}

\begin{table}[H]
    \centering
    \caption{Logiciels requis}
    \begin{tabular}{|l|l|l|}
        \hline
        \textbf{Logiciel} & \textbf{Version min} & \textbf{Lien} \\
        \hline
        Azure CLI & 2.40+ & microsoft.com/cli \\
        Terraform & 1.0+ & terraform.io \\
        Ansible & 2.10+ & ansible.com \\
        Docker & 20.10+ & docker.com \\
        kubectl & 1.20+ & kubernetes.io \\
        Git & 2.30+ & git-scm.com \\
        \hline
    \end{tabular}
\end{table}

\subsection{Comptes requis}

\begin{itemize}
    \item Compte Microsoft Azure (Free Tier accepté)
    \item Compte GitHub
    \item Compte Docker Hub
\end{itemize}

\section{Étape 1: Préparation de l'environnement}

\subsection{1.1 Authentification Azure}

\begin{lstlisting}[language=bash]
# Login
az login

# Configurer la souscription par défaut
az account set --subscription "YOUR_SUBSCRIPTION_ID"

# Vérifier
az account show
\end{lstlisting}

\subsection{1.2 Clonage du repository}

\begin{lstlisting}[language=bash]
git clone https://github.com/username/devops-project.git
cd devops-project
\end{lstlisting}

\subsection{1.3 Configuration Docker}

\begin{lstlisting}[language=bash]
docker login
# Entrez vos credentials Docker Hub
\end{lstlisting}

\section{Étape 2: Déploiement de l'infrastructure avec Terraform}

\subsection{2.1 Initialisation Terraform}

\begin{lstlisting}[language=bash]
cd terraform
terraform init
\end{lstlisting}

Terraform télécharge les providers requis.

\subsection{2.2 Planification}

\begin{lstlisting}[language=bash]
terraform plan
\end{lstlisting}

Affiche les ressources qui seront créées.

\subsection{2.3 Application}

\begin{lstlisting}[language=bash]
terraform apply --auto-approve
\end{lstlisting}

\textbf{Temps estimé :} 10-15 minutes.

\subsection{2.4 Récupération des IPs}

\begin{lstlisting}[language=bash]
terraform output
\end{lstlisting}

Les IPs publiques des VMs seront affichées.

\section{Étape 3: Configuration avec Ansible}

\subsection{3.1 Permissions SSH}

\begin{lstlisting}[language=bash]
cd ../ansible
chmod 600 ../terraform/id_rsa_devops
\end{lstlisting}

\subsection{3.2 Désactivation de la vérification SSH}

\begin{lstlisting}[language=bash]
export ANSIBLE_HOST_KEY_CHECKING=False
\end{lstlisting}

\subsection{3.3 Exécution du playbook}

\begin{lstlisting}[language=bash]
ansible-playbook -i inventory.ini playbook.yml
\end{lstlisting}

\textbf{Temps estimé :} 10-15 minutes.

\section{Étape 4: Configuration de kubectl}

\subsection{4.1 Copie du kubeconfig}

Le fichier kubeconfig est automatiquement récupéré par Terraform et Ansible:

\begin{lstlisting}[language=bash]
# Déplacer vers le répertoire kubectl standard
mkdir -p ~/.kube
cp kubeconfig.yaml ~/.kube/config
\end{lstlisting}

\subsection{4.2 Verification de la connexion}

\begin{lstlisting}[language=bash]
kubectl cluster-info

# Doit afficher quelque chose comme:
# Kubernetes control-plane is running at https://...
\end{lstlisting}

\subsection{4.3 Vérification des nodes}

\begin{lstlisting}[language=bash]
kubectl get nodes

# Doit afficher:
# NAME          STATUS   ROLES                  READY
# <master-ip>   Ready    control-plane,master   1/1
# <worker-ip>   Ready    <none>                 1/1
\end{lstlisting}

\section{Étape 5: Construction de l'image Docker}

\subsection{5.1 Build local}

\begin{lstlisting}[language=bash]
docker build -t task-manager:latest .
\end{lstlisting}

\subsection{5.2 Tagging}

\begin{lstlisting}[language=bash]
docker tag task-manager:latest \
  seiflegnioui/task-manager:latest
\end{lstlisting}

\subsection{5.3 Push vers Docker Hub}

\begin{lstlisting}[language=bash]
docker push seiflegnioui/task-manager:latest
\end{lstlisting}

\section{Étape 6: Déploiement sur Kubernetes}

\subsection{6.1 Application des manifests}

\begin{lstlisting}[language=bash]
kubectl apply -f kubernetes/
\end{lstlisting}

\subsection{6.2 Vérification du déploiement}

\begin{lstlisting}[language=bash]
# Attendre que les pods soient Ready
kubectl get pods -w

# Ctrl+C pour quitter
\end{lstlisting}

\subsection{6.3 Vérification du service}

\begin{lstlisting}[language=bash]
kubectl get svc

# Doit afficher task-manager-service sur le port 30080
\end{lstlisting}

\section{Étape 7: Accès à l'application}

\subsection{7.1 Déterminer l'URL}

\begin{lstlisting}[language=bash]
# Récupérer l'IP du master
MASTER_IP=\$(terraform output -raw master_public_ip)
echo "Application accessible à: http://\$MASTER_IP:30080"
\end{lstlisting}

\subsection{7.2 Test de l'API}

\begin{lstlisting}[language=bash]
# Récupérer les tâches
curl http://\$MASTER_IP:30080/tasks

# Créer une tâche
curl -X POST http://\$MASTER_IP:30080/tasks \
  -H "Content-Type: application/json" \
  -d '{\"title\": \"Test Task\", \"description\": \"Description\"}'

# Vérifier la santé
curl http://\$MASTER_IP:30080/health
\end{lstlisting}

\section{Étape 8: Configuration du pipeline CI/CD}

\subsection{8.1 Secrets GitHub}

Sur GitHub, aller à Settings $\rightarrow$ Secrets and variables $\rightarrow$ Actions:

\begin{enumerate}
    \item \textbf{DOCKER\_USERNAME} : Votre username Docker Hub
    \item \textbf{DOCKER\_PASSWORD} : Votre token Docker Hub
    \item \textbf{KUBECONFIG} : Contenu base64 de kubeconfig.yaml
\end{enumerate}

Générer la version base64:

\begin{lstlisting}[language=bash]
cat kubeconfig.yaml | base64 -w 0
\end{lstlisting}

\subsection{8.2 Activation du workflow}

Créer le fichier \texttt{.github/workflows/deploy.yml} avec la configuration du workflow.

\subsection{8.3 Test du workflow}

\begin{lstlisting}[language=bash]
git push origin main
\end{lstlisting}

Le workflow s'exécute automatiquement.

\section{Troubleshooting}

\subsection{Erreurs de connexion Terraform}

\begin{lstlisting}[language=bash]
# Vérifier l'authentification Azure
az account show

# Réinitialiser Terraform
rm -rf .terraform .terraform.lock.hcl
terraform init
\end{lstlisting}

\subsection{Erreurs SSH Ansible}

\begin{lstlisting}[language=bash]
# Tester la connectivité SSH
ssh -i terraform/id_rsa_devops azureuser@<master-ip>

# Vérifier les permissions
chmod 600 terraform/id_rsa_devops
\end{lstlisting}

\subsection{Pods non Ready}

\begin{lstlisting}[language=bash]
# Voir les événements
kubectl describe pod <pod-name>

# Voir les logs
kubectl logs <pod-name>

# Recréer les pods
kubectl rollout restart deployment/task-manager-api
\end{lstlisting}

\newpage

%% Chapter 10: Maintenance et Monitoring
\chapter{Maintenance, Monitoring et Meilleures Pratiques}

\section{Monitoring des applications}

\subsection{Endpoins de santé}

L'application expose un endpoint health check:

\begin{lstlisting}[language=bash]
curl http://master-ip:30080/health
# Réponse: {"status": "healthy"}
\end{lstlisting}

\subsection{Kubernetes health checks}

\subsubsection{Readiness Probe}

Configurée pour vérifier la disponibilité de l'application:

\begin{lstlisting}[language=yaml]
readinessProbe:
  httpGet:
    path: /health
    port: 5000
  initialDelaySeconds: 10
  periodSeconds: 5
\end{lstlisting}

\subsubsection{Liveness Probe}

À ajouter au deployment pour l'auto-healing:

\begin{lstlisting}[language=yaml]
livenessProbe:
  httpGet:
    path: /health
    port: 5000
  initialDelaySeconds: 30
  periodSeconds: 10
  failureThreshold: 3
\end{lstlisting}

\section{Logging et debugging}

\subsection{Logs Kubernetes}

\begin{lstlisting}[language=bash]
# Logs actuels
kubectl logs <pod-name>

# Logs en temps réel
kubectl logs -f <pod-name>

# Logs pour tous les pods du déploiement
kubectl logs -l app=task-manager

# Logs du pod précédent (si crash)
kubectl logs <pod-name> --previous
\end{lstlisting}

\subsection{Inspection des ressources}

\begin{lstlisting}[language=bash]
# Détails Pod
kubectl describe pod <pod-name>

# Détails Deployment
kubectl describe deployment task-manager-api

# Détails Service
kubectl describe service task-manager-service

# Événements du cluster
kubectl get events
\end{lstlisting}

\section{Gestion des versions}

\subsection{Versioning d'images}

Utiliser le semantic versioning (MAJOR.MINOR.PATCH):

\begin{lstlisting}[language=bash]
# Tags multiples
docker tag task-manager:vv1.0.0 seiflegnioui/task-manager:v1.0.0
docker tag task-manager:v1.0.0 seiflegnioui/task-manager:latest
docker push seiflegnioui/task-manager:v1.0.0
docker push seiflegnioui/task-manager:latest
\end{lstlisting}

\subsection{Déploiement de nouvelles versions}

\begin{lstlisting}[language=bash]
# Mise à jour du déploiement
kubectl set image deployment/task-manager-api \
  task-manager-api=seiflegnioui/task-manager:v2.0.0

# Ou avec apply + YAML modifié
kubectl apply -f kubernetes/deployment.yaml
\end{lstlisting}

\section{Meilleures pratiques}

\subsection{Infrastructure}

\begin{enumerate}
    \item Utiliser des tagsimages spécifiques (jamais "latest" en prod)
    \item Implémenter des Network Policies pour la sécurité
    \item Utiliser des Storage Classes pour la persistance
    \item Configurer des Resource Quotas par namespace
    \item Utiliser RBAC (Role-Based Access Control)
\end{enumerate}

\subsection{Application}

\begin{enumerate}
    \item Implémenter des health checks
    \item Utiliser des probes de liveness et readiness
    \item Gérer les signaux SIGTERM pour un arrêt gracieux
    \item Implémenter le circuit breaker et retry logic
    \item Utiliser des logs structurés (JSON)
\end{enumerate}

\subsection{DevOps}

\begin{enumerate}
    \item Versioner toute l'infrastructure (IaC)
    \item Implémenter des tests en CI/CD
    \item Utiliser des semantic versioning pour les releases
    \item Sauvegarder les états Terraform
    \item Documenter les runbooks
\end{enumerate}

\subsection{Sécurité}

\begin{enumerate}
    \item Scannner les images avec Trivy/Aqua
    \item Utiliser des namespaces Kubernetes
    \item Implémenter Network Policies
    \item Chiffrer les secrets (Sealed Secrets/Vault)
    \item Auditerance les accès avec RBAC
\end{enumerate}

\section{Cleanup et destruction}

\subsection{Destruction de l'infrastructure}

\begin{lstlisting}[language=bash]
# Depuis le répertoire terraform
terraform destroy --auto-approve
\end{lstlisting}

Cette commande supprime:
\begin{itemize}
    \item Les VMs Azure
    \item Le Virtual Network
    \item Le Network Security Group
    \item Les adresses IP publiques
    \item Tous les disques associés
\end{itemize}

\textbf{Attention :} Cette opération supprime définitivement tous les résources.

\newpage

%% Chapter 11: Conclusions
\chapter{Conclusions et Perspectives}

\section{Résumé du projet}

Ce projet a démontré avec succès l'implémentation d'une pile DevOps complète et moderne:

\begin{itemize}
    \item Déploiement d'une application Flask REST
    \item Containerisation avec Docker
    \item Infrastructure as Code avec Terraform
    \item Configuration automation avec Ansible
    \item Orchestration Kubernetes (K3s)
    \item Pipeline CI/CD automatisé
    \item Documentation complète et guides
\end{itemize}

\section{Technologies et outils maîtrisés}

\begin{enumerate}
    \item \textbf{Langage} : Python (Flask, Gunicorn)
    \item \textbf{Containerisation} : Docker
    \item \textbf{Cloud} : Microsoft Azure
    \item \textbf{IaC} : Terraform
    \item \textbf{Configuration} : Ansible
    \item \textbf{Orchestration} : Kubernetes (K3s)
    \item \textbf{CI/CD} : GitHub Actions
    \item \textbf{Contrôle de version} : Git
\end{enumerate}

\section{Améliorations futures}

\subsection{Court terme}

\begin{itemize}
    \item Ajouter des tests unitaires et d'intégration
    \item Implémenter un monitoring avec Prometheus/Grafana
    \item Ajouter une base de données persistante
    \item Configurer TLS/SSL pour la sécurité
    \item Implémenter le logging avec ELK Stack
\end{itemize}

\subsection{Moyen terme}

\begin{itemize}
    \item Migrer vers AKS (Azure Kubernetes Service)
    \item Implémenter Helm pour la gestion des charts
    \item Ajouter ArgoCD pour GitOps
    \item Configurer des Network Policies
    \item Implémenter Service Mesh (Istio/Linkerd)
\end{itemize}

\subsection{Long terme}

\begin{itemize}
    \item Migration vers le multi-cloud
    \item Disaster recovery et backup automatisé
    \item Implémentation de chaos engineering
    \item ML Ops pour les modèles de ML
    \item Serverless pour les opérations ponctuelles
\end{itemize}

\section{Leçons apprises}

\subsection{Points clés}

\begin{enumerate}
    \item L'automatisation est essentielle en DevOps
    \item Infrastructure as Code rend l'infrastructure versionnable
    \item La conteneurisation améliore la reproductibilité
    \item L'orchestration Kubernetes est puissante mais complexe
    \item Les pipelines CI/CD accélèrent très significativement les déploiements
    \item La documentation est cruciale pour la maintenance
\end{enumerate}

\subsection{Défis rencontrés}

\begin{itemize}
    \item Configuration SSH avec Ansible sur Azure
    \item Gestion des IPs dynamiques et templates
    \item Synchronisation des tokens dans l'orchestration
    \item Configuration des probes de santé appropriées
\end{itemize}

\section{Recommandations}

Pour toute organisation déployant une architecture similaire:

\begin{enumerate}
    \item Commencer simple et itérer progressivement
    \item Automatiser depuis le début
    \item Investir dans la documentation
    \item Utiliser des standards de l'industrie
    \item Implémenter la monitoring et les alertes
    \item Faire des tests de disaster recovery régulièrement
    \item Maintenir une culture DevOps forte
\end{enumerate}

\newpage

%% Appendices
\appendix

\chapter{Code Source Complet}

\section{Application Flask (app.py)}

\begin{lstlisting}[language=Python]
from flask import Flask, jsonify, request, render_template
import uuid

app = Flask(__name__)

tasks = []

@app.route('/')
def index():
    return render_template('index.html')

@app.route('/tasks', methods=['GET'])
def get_tasks():
    return jsonify(tasks), 200

@app.route('/tasks', methods=['POST'])
def create_task():
    data = request.get_json()
    if not data or not 'title' in data:
        return jsonify({'error': 'Title is required'}), 400
    
    task = {
        'id': str(uuid.uuid4()),
        'title': data['title'],
        'description': data.get('description', ''),
        'done': False
    }
    tasks.append(task)
    return jsonify(task), 201

@app.route('/tasks/<task_id>', methods=['PATCH', 'PUT'])
def update_task(task_id):
    data = request.get_json()
    for task in tasks:
        if task['id'] == task_id:
            if 'done' in data:
                task['done'] = data['done']
            return jsonify(task), 200
    return jsonify({'error': 'Task not found'}), 404

@app.route('/tasks/<task_id>', methods=['DELETE'])
def delete_task(task_id):
    global tasks
    tasks = [task for task in tasks if task['id'] != task_id]
    return jsonify({'result': True}), 200

@app.route('/health', methods=['GET'])
def health_check():
    return jsonify({'status': 'healthy'}), 200

if __name__ == '__main__':
    app.run(host='0.0.0.0', port=5000, debug=False)
\end{lstlisting}

\chapter{Configuration Terraform}

\section{Fichier main.tf complet}

Voir le fichier \texttt{terraform/main.tf} dans le repository pour la configuration complète des ressources Azure.

\chapter{Manifests Kubernetes}

\section{Déploiement complet}

Voir le répertoire \texttt{kubernetes/} pour tous les manifests YAML.

\chapter{Résumé des commandes}

\section{Commandes essentielles}

\begin{longtable}{|p{4cm}|p{5cm}|}
    \hline
    \textbf{Tâche} & \textbf{Commande} \\
    \hline
    \endhead
    \hline
    \multicolumn{2}{r}{(Continued on next page)} \\
    \endfoot
    \hline
    \endlastfoot
    Init Terraform & \texttt{terraform init} \\
    Plan Terraform & \texttt{terraform plan} \\
    Apply Terraform & \texttt{terraform apply} \\
    Destroy infra & \texttt{terraform destroy} \\
    Run Ansible & \texttt{ansible-playbook -i inventory.ini playbook.yml} \\
    Build Docker & \texttt{docker build -t name:tag .} \\
    Push Docker & \texttt{docker push name:tag} \\
    Deploy K8s & \texttt{kubectl apply -f kubernetes/} \\
    Get pods & \texttt{kubectl get pods} \\
    Get logs & \texttt{kubectl logs pod-name} \\
    Describe pod & \texttt{kubectl describe pod pod-name} \\
    Port forward & \texttt{kubectl port-forward pod-name 5000:5000} \\
    Scale deployment & \texttt{kubectl scale deployment name --replicas=3} \\
    Restart deployment & \texttt{kubectl rollout restart deployment/name} \\
\end{longtable}

\newpage

\chapter{Ressources et références}

\section{Documentation officielle}

\begin{itemize}
    \item Terraform Documentation : \url{https://www.terraform.io/docs}
    \item Kubernetes Documentation : \url{https://kubernetes.io/docs}
    \item Ansible Documentation : \url{https://docs.ansible.com}
    \item Docker Documentation : \url{https://docs.docker.com}
    \item Azure Documentation : \url{https://docs.microsoft.com/azure}
    \item GitHub Actions : \url{https://docs.github.com/en/actions}
\end{itemize}

\section{Livres recommandés}

\begin{itemize}
    \item "The DevOps Handbook" - Gene Kim, Patrick Debois
    \item "Infrastructure as Code" - Kief Morris
    \item "Kubernetes in Action" - Marko Luksa
    \item "The Phoenix Project" - Gene Kim, Kevin Behr
\end{itemize}

\section{Community et forums}

\begin{itemize}
    \item Stack Overflow : \url{https://stackoverflow.com/}
    \item Reddit r/devops : \url{https://reddit.com/r/devops/}
    \item CNCF Slack : \url{https://cloud-native.slack.com/}
\end{itemize}

\newpage

%% Final
\chapter*{Conclusion Générale}
\addcontentsline{toc}{chapter}{Conclusion Générale}

Ce projet représente une démonstration complète des principes et outils DevOps modernes. À travers la conception, l'implémentation et le déploiement d'une architecture cloud entièrement automatisée, nous avons acquis une expérience pratique des technologies les plus exigeantes du secteur.

\vspace{1cm}

L'architecture proposée est scalable, sécurisée et durable, capable de supporter une croissance future et des modifications. Les pipelines CI/CD automatisés garantissent une livraison rapide et fiable des mises à jour logicielles.

\vspace{1cm}

Les apprentissages de ce projet s'étendent bien au-delà des outils spécifiques utilisés - ils englobent les principes fondamentaux du DevOps: l'automatisation, la collaboration, la mesure et le partage. Ces principes resteront pertinents peu importe l'évolution technologique.

\vspace{1cm}

Nous sommes convaincus que cette fondation solide permettra une croissance future et una adaptation aux besoins changeants de l'industrie.

\newpage

%% Document Information
\begin{flushright}
    \large
    Rapport généré le \today\\
    Module MEIA - Cloud DevOps Engineering\\
    Version 1.0
\end{flushright}

\end{document}
